기룡이는 한 주에 평일 $5$일씩 $4$주간 $5$천 원짜리 아이스 카페 라떼를 마시면 한 달에 무려 $10$만 원이란 사실을 깨달았다. 기룡이는 용돈을 아끼기 위해 가계부를 쓰기 시작했고, 한 달 동안 어디에 돈을 가장 많이 썼는지 알아보기 위해 가계부를 정산하기로 했다.

기룡이의 가계부에는 지출 날짜, 카테고리, 지출 금액이 기록되어 있다. 카테고리는 음식, 취미생활, 생활비, 의료비, 교통비의 $5$가지로 구분된다.

한 달 동안 기록한 지출 내역이 주어질 때, 카테고리별 총지출 금액을 계산하여 총지출 금액이 가장 큰 카테고리와 가장 작은 카테고리, 그리고 각각의 총지출 금액을 구해보자.
